\medskip

Une équipe de biologistes étudie l'évolution de la superficie recouverte par une algue marine appelée posidonie, sur le fond de la baie de l'Alycastre, près de l'île de Porquerolles.

La zone étudiée est d'une superficie totale de 20 hectares (ha), et au premier juillet 2024, la posidonie recouvrait 1 ha de cette zone.

\begin{center}\textbf{Partie A :} étude d'un modèle discret\end{center}

Pour tout entier naturel $n$, on note $u_n$ la superficie de la zone, en hectare, recouverte
par la posidonie au premier juillet de l'année $2024 + n$. Ainsi, $u_0 = 1$.

\smallskip

Une étude conduite sur cette superficie a permis d'établir que pour tout entier
naturel $n$ :

\[u_{n+1} = - 0,02u_n^2 + 1,3u_n.\]

\smallskip

\begin{enumerate}
\item Calculer la superficie que devrait recouvrir la posidonie au premier juillet 2025 d'après ce modèle.
\item On note $h$ la fonction définie sur [0~;~20] par 

\[h(x) = - 0,02x^2 + 1,3x.\]

On admet que $h$ est croissante sur [0~;~20].
	\begin{enumerate}
		\item Démontrer que pour tout entier naturel $n,\: 1 \leqslant u_n  \leqslant u_{n+1} \leqslant 20$.
		\item En déduire que la suite $\left(u_n\right)$ converge. On note $L$ sa limite.
		\item Justifier que $L = 15$.
	\end{enumerate}
\end{enumerate}
\begin{minipage}{0.65\linewidth}
\begin{enumerate}[resume]
\item Les biologistes souhaitent savoir au bout de combien de temps la surface recouverte par la
posidonie dépassera les 14 hectares.
	\begin{enumerate}
		\item Sans aucun calcul, justifier que, d'après ce modèle, cela se produira.
		\item Recopier et compléter l'algorithme suivant pour qu'en fin d'exécution, il affiche la réponse à la question des biologistes.
	\end{enumerate}
\end{enumerate}
\end{minipage}\hfill
\begin{minipage}{0.3\linewidth}
\begin{Verbatim}[frame=single]
def seuil():
  n=0
  u= 1
  while ...... :
    n=......
    u=......
  return n
\end{Verbatim}
\end{minipage}

\begin{center} \textbf{Partie B :} étude d'un modèle continu\end{center}

On souhaite décrire la superficie de la zone étudiée recouverte par la posidonie au cours du temps avec un modèle continu.

Dans ce modèle, pour une durée $t$, en année, écoulée à partir du premier juillet 2024, la superficie de la zone étudiée recouverte par la posidonie est donnée par
$f(t)$, où $f$ est une fonction définie sur $[0~;~ +\infty[$ vérifiant :

\begin{itemize}[label=$\bullet~$]
\item $f(0) = 1$ ;
\item $f$ ne s'annule pas sur $[0~;~ +\infty[$ ;
\item $f$ est dérivable sur $[0~;~ +\infty[$ ;
\item $f$ est solution sur $[0~;~ +\infty[$ de l'équation différentielle 

\[(E_1) : \quad y' =0,02y(15 - y).\]

\end{itemize}

On admet qu'une telle fonction $f$ existe; le but de cette partie est d'en déterminer une expression.

On note $f'$ la fonction dérivée de $f$.

\medskip

\begin{enumerate}
\item Soit $g$ la fonction définie sur $[0~;~ +\infty[$ par $g(t) = \dfrac{1}{f(t)}$.

Montrer que $g$ est solution de l'équation différentielle 

\[(E_2) : \quad y' = - 0,3y + 0,02.\]

\item Donner les solutions de l'équation différentielle $(E_2)$.
\item En déduire que pour tout $t \in  [0~;~ +\infty[$ :

\[f(t)=  \dfrac{15}{14\e^{-0,3t} + 1}.\]

\item Déterminer la limite de $f$ en $+\infty$.
\item Résoudre dans l'intervalle $[0~;~ +\infty[$ l'inéquation $f(t) > 14$. Interpréter le résultat dans le contexte de l'exercice.
\end{enumerate}
